% Inline pipeline figure: an input is packaged into an artifact, a loop runs
% against it under a fixed budget, and a held-out checker turns the result into
% a score. Drop this file into your tex/ and \input it inside the document.
%
% Preamble the figure needs:
%   \usepackage{tikz}
%   \usetikzlibrary{arrows,calc,positioning}
%   \usepackage{caption}                    % for \captionof outside a float
%   \definecolor{ggrey}{HTML}{5F6368}
%   \definecolor{gblue}{HTML}{4285F4}
%   \definecolor{gred}{HTML}{EA4335}
%   \definecolor{ggreen}{HTML}{34A853}
%
% Colour follows DOCTRINE.md rule 4: blue is the step the method owns, red is
% what is withheld from it, green is what comes out, grey is neutral. The flows
% are grey throughout, because an arrow is a relation rather than a party.
\begin{figure}[h]
\centering
\begin{tikzpicture}[
  font=\sffamily\fontsize{8}{10}\selectfont,
  box/.style={draw=ggrey!45, line width=0.5pt, rounded corners=1pt, align=center,
              inner xsep=5pt, inner ysep=5pt, text width=62pt},
  owned/.style={box, draw=gblue!60, fill=gblue!6},
  withheld/.style={box, draw=gred!55, fill=gred!6},
  result/.style={box, draw=ggreen!60, fill=ggreen!7},   % not 'out': tikz reserves it
  lbl/.style={font=\sffamily\fontsize{6.6}{8}\selectfont, color=ggrey, align=center},
  flow/.style={-latex, draw=ggrey!70, line width=0.6pt},
]

\node[box]      (src)   at (0,0)     {Input artifact\\[2pt]
                                      {\fontsize{6.6}{8}\selectfont what already exists:\\ code, data, a paper}};
\node[owned]    (pack)  at (2.75,0)  {Packaging step\\[2pt]
                                      {\fontsize{6.6}{8}\selectfont the part the method owns}};
\node[box]      (env)   at (5.50,0)   {Runnable target\\[2pt]
                                      {\fontsize{6.6}{8}\selectfont harness, budget,\\ submission API}};
\node[result]   (index) at (11.35,0)  {Result record\\[2pt]
                                      {\fontsize{6.6}{8}\selectfont score, cost,\\ full trajectory}};

\node[box]      (loop)  at (8.60,1.05)  {Agent loop\\[2pt]
                                         {\fontsize{6.6}{8}\selectfont propose, run,\\ read, revise}};
\node[withheld] (check) at (8.60,-1.15) {Hidden checker\\[2pt]
                                         {\fontsize{6.6}{8}\selectfont frozen reference,\\ held-out metric}};

\draw[flow] (src)  -- (pack);
\draw[flow] (pack) -- (env);
\draw[flow] (env.east) -- ++(0.30,0) |- (loop.west);
\draw[flow] (loop.east) -| ($(check.north)+(0.75,0)$) -- ($(check.north east)+(-0.2,0)$);
\draw[flow] (check.east) -- ++(0.30,0) -| (index.south);

% Small caps riders name the edge relation without adding a legend.
\node[lbl, above=1pt of src]   {\textsc{input}};
\node[lbl, above=1pt of pack]  {\textsc{packaged as}};
\node[lbl, above=1pt of index] {\textsc{producing}};
\node[lbl, anchor=west] at ($(loop.east)+(0.12,0.42)$)  {submission};
\node[lbl, anchor=east] at (6.95,1.05) {fixed budget};

\end{tikzpicture}
\vspace{0.4em}
\captionof{figure}{\textbf{From an input artifact to a scored result.} State in the caption what
each box is, what the agent may touch, and what it never sees. A caption that names the withheld
component is what makes the red fill readable without a legend.}
\label{fig:pipeline}
\vspace{0.6em}
\end{figure}
